\section{Conclusion}

Cross-network traffic-signal transfer improved when the learned object was
reduced from absolute dynamics to normalized, matched action contrasts.
Anchoring a parent-restricted score and correcting it with an antisymmetric
all-action model reduced seed-held-out action regret in development and in two
external city-derived networks.  The same frozen models improved closed-loop
waiting time over a dense residual and produced favourable simulator-only and
rigid-anchor point estimates, while remaining behind MaxPressure and phase
pressure.  Crucially, a 64-seed target-label-only comparison rejected a general
source-contribution claim: source-plus-target fitting improved Jinan but harmed
Los Angeles and was worse in the equal-city mean.  CFCMT should therefore be
interpreted as an auditable action-contrast architecture, not evidence that
source data are uniformly beneficial.  Establishing cross-city source benefit
will require a precommitted source-selection rule, more independent target
cities and either logged interventions, a validated digital twin or prospective
field control; the present SUMO counterfactuals do not supply that evidence.
